\documentclass[../../../include/open-logic-section]{subfiles}

\begin{document}
\olfileid{sth}{card-arithmetic}{simp}

\olsection{સરવાળો અને ગુણાકારનું સરળીકરણ}

અનંતાતીત ગણાંક-સરવાળો અને ગણાંક-ગુણાકાર
\emph{અત્યંત} સહેલા નીવડે છે. કારણ કે ગણાંકો
ચોક્કસ ક્રમવાચક સંખ્યાઓ છે; તેથી તેઓ સુક્રમિત
છે અને એક ચોક્કસ રીતે તેમનો ઉપયોગ કરી શકાય છે.
જોકે આ બતાવવું \emph{એટલું} સહેલું નથી. શરૂઆતમાં
થોડી ચતુર વ્યાખ્યા જોઈએ:

\begin{defn}
ક્રમવાચક સંખ્યાઓની જોડીઓ પર \emph{પ્રમાણભૂત ક્રમ}
$\canonord$ આમ વ્યાખ્યાયિત કરીએ:
$\tuple{\alpha_1, \alpha_2} \canonord
\tuple{\beta_1, \beta_2}$ જો અને માત્ર જો નીચેનામાંથી
કોઈ એક શરત સાચી હોય:
\begin{enumerate}
	\item $\max(\alpha_1, \alpha_2) < \max(\beta_1, \beta_2)$; અથવા
	\item $\max(\alpha_1, \alpha_2) = \max(\beta_1, \beta_2)$ અને
	$\alpha_1 < \beta_1$; અથવા
	\item $\max(\alpha_1, \alpha_2) = \max(\beta_1, \beta_2)$ અને
	$\alpha_1 = \beta_1$ અને $\alpha_2 < \beta_2$
\end{enumerate}
\end{defn}

\begin{lem}
કોઈ પણ ક્રમવાચક સંખ્યા $\alpha$ માટે
$\tuple{\alpha \times \alpha, \canonord}$ સુક્રમ છે.
\end{lem}

\begin{proof}
$\alpha \times \alpha$ પર $\canonord$ પૂર્ણ તુલનાત્મક છે
તે સ્પષ્ટ છે. કારણ કે ધારો કે
$\tuple{\alpha_1, \alpha_2}$ અને
$\tuple{\beta_1, \beta_2}$ પૈકી કોઈ પણ બીજી કરતાં
$\canonord$ મુજબ નાની નથી. તો
$\max(\alpha_1, \alpha_2) = \max(\beta_1, \beta_2)$
અને $\alpha_1 = \beta_1$ તથા $\alpha_2 = \beta_2$;
તેથી $\tuple{\alpha_1, \alpha_2} =
\tuple{\beta_1,  \beta_2}$.

સુક્રમિતતા બતાવવા અરિક્ત
$X \subseteq \alpha\times\alpha$ લો. $\alpha$
ક્રમવાચક સંખ્યા હોવાથી
$\Setabs{\max(\gamma_1, \gamma_2)}{\tuple{\gamma_1,
\gamma_2} \in X}$માં કોઈ લઘુતમ ઘટક $\delta$ છે.
હવે $\Setabs{\tuple{\gamma_1,\gamma_2} \in X}{\max(\gamma_1, \gamma_2) =
\delta}$માંથી લઘુતમ પ્રથમ ઘટક ન ધરાવતી બધી જોડીઓ
દૂર કરો. બાકીની જોડીઓમાં લઘુતમ બીજો ઘટક ધરાવતી
જોડી $X$ની $\canonord$ મુજબ લઘુતમ જોડી છે.
\end{proof}
\noindent
હવે એક નાનું, સરળ અવલોકન:

\begin{prop}\ollabel{simplecardproduct}
જો $\cardeq{\alpha}{\beta}$ હોય, તો
$\cardeq{\alpha \times \alpha}{\beta \times \beta}$.
\end{prop}

\begin{proof}
$f \colon \alpha \to \beta$ પરથી
$\tuple{\gamma_1, \gamma_2} \mapsto
\tuple{f(\gamma_1), f(\gamma_2)}$ લો.
\end{proof}
\noindent
હવે આ બધું એક નિર્ણાયક લેમાના પુરાવામાં વાપરીશું:
\begin{lem}\ollabel{alphatimesalpha}
કોઈ પણ અનંત ક્રમવાચક સંખ્યા $\alpha$ માટે
$\cardeq{\alpha}{\alpha \times \alpha}$.
\end{lem}

\begin{proof}
વિરોધાભાસ માટે ધારો કે આ ખોટું પડે એવી
લઘુતમ અનંત ક્રમવાચક સંખ્યા $\alpha$ છે.
\olref[sfr][siz][zigzag]{natsquaredenumerable}
બતાવે છે કે $\cardeq{\omega}{\omega\times\omega}$;
તેથી $\omega \in \alpha$. વધુમાં $\alpha$ ગણાંક છે:
વિરોધાભાસ માટે અન્યથા ધારો. તો
$\card{\alpha} \in \alpha$, તેથી અનુમાન મુજબ
$\cardeq{\card{\alpha}}{\card{\alpha} \times
\card{\alpha}}$; અને વ્યાખ્યાથી
$\cardeq{\card{\alpha}}{\alpha}$; તેથી
\olref{simplecardproduct}થી
$\cardeq{\alpha}{\alpha\times\alpha}$.

હવે દરેક $\tuple{\gamma_1, \gamma_2} \in
\alpha \times \alpha$ માટે આ ખંડ જુઓ:
\begin{align*}
	\text{Seg}(\gamma_1, \gamma_2) &= \Setabs{\tuple{\delta_1, \delta_2} \in \alpha \times \alpha}{\tuple{\delta_1, \delta_2} \canonord \tuple{\gamma_1, \gamma_2}}
\end{align*}
$\gamma = \max(\gamma_1, \gamma_2)$ લો. નોંધો કે
$\tuple{\gamma_1, \gamma_2} \canonord
\tuple{\gamma+1, \gamma + 1}$. તેથી $\gamma$
અનંત હોય ત્યારે:
\begin{align*}
	\text{Seg}(\gamma_1, \gamma_2) & 
	\precsim ((\gamma \ordplus 1)\ordtimes (\gamma \ordplus 1))\\
	&\approx (\gamma \ordtimes \gamma)
	\text{, \olref[ord-arithmetic][using-addition]{ordinfinitycharacter} અને 
	\olref{simplecardproduct} મુજબ}\\
	&\approx \gamma \text{, અનુમાનની ધારણા મુજબ}\\
	& \prec \alpha\text{, કારણ કે $\alpha$ ગણાંક છે}
\end{align*}
આથી $\ordtype{\alpha\times \alpha, \canonord} \leq
\alpha$, અને તેથી $\cardle{\alpha \times
\alpha}{\alpha}$. નિશ્ચિતપણે
$\cardle{\alpha}{\alpha \times \alpha}$ પણ છે;
તેથી શ્રેડર--બર્નસ્ટાઇનથી પરિણામ મળે છે.
\end{proof}

અંતે આપણને ઇચ્છેલું સરળીકરણ મળે છે:

\begin{thm}\ollabel{cardplustimesmax}
$\cardfont{a}, \cardfont{b}$ અનંત ગણાંકો હોય, તો:
\[
	\cardfont{a}
\cardtimes \cardfont{b} = \cardfont{a} \cardplus \cardfont{b} =
\text{max}(\cardfont{a}, \cardfont{b}).
\]
\end{thm}

\begin{proof}
સામાન્યતા ગુમાવ્યા વિના
$\cardfont{a} = \max(\cardfont{a},
\cardfont{b})$ ધારો. પછી \olref{alphatimesalpha}
વાપરીને
$\cardfont{a}\cardtimes\cardfont{a} = \cardfont{a} \leq \cardfont{a}
\cardplus \cardfont{b} \leq \cardfont{a} \cardplus \cardfont{a} \leq
\cardfont{a} \cardtimes \cardfont{a}$.
\end{proof}\noindent એ જ રીતે
$\cardfont{a}$ અનંત હોય, તો કદ
$\leq\cardfont{a}$ ધરાવતા ગણોના $\cardfont{a}$
કદના યોગગણનું કદ પણ $\leq\cardfont{a}$ છે:

\begin{prop}\ollabel{kappaunionkappasize}
ધારો કે $\cardfont{a}$ અનંત ગણાંક છે. દરેક
ક્રમવાચક સંખ્યા $\beta \in \cardfont{a}$ માટે
$X_\beta$ એવો ગણ લો કે
$\card{X_\beta} \leq \cardfont{a}$. તો
$\card{\bigcup_{\beta \in \cardfont{a}} X_\beta}
\leq \cardfont{a}$.
\end{prop}

\begin{proof}
દરેક $\beta \in \cardfont{a}$ માટે
!!a{injection} $f_\beta \colon X_\beta \to
\cardfont{a}$ નક્કી કરો.\footnote{આ વિધેયો
``નક્કી'' કેવી રીતે કરાય છે? જુઓ
\olref[sth][choice][countablechoice]{sec}.}
!!a{injection} $g \colon \bigcup_{\beta \in
\cardfont{a}} X_\beta \to \cardfont{a} \times
\cardfont{a}$ને $g(v) = \tuple{\beta, f_\beta(v)}$
વડે વ્યાખ્યાયિત કરો, જ્યાં $v \in X_\beta$
અને દરેક $\gamma \in \beta$ માટે
$v \notin X_\gamma$. હવે
\olref{cardplustimesmax} મુજબ
$\bigcup_{\beta \in \cardfont{a}} X_\beta \preceq
\cardfont{a} \times \cardfont{a} \approx \cardfont{a}$.
\end{proof}

\end{document}
